% Section 1.8, from lectures/L01/appendix/b_exercises.md.

\section{Exercises}
\label{c1:sec:exercises}

\begin{aside}
\textbf{How to work these.} Exercise~\ref{c1:ex:1} is a short pen-and-paper warm-up on notation this
book uses throughout; it needs no VHDL, GHDL, or board. Each exercise is tagged with the kinds of
work it asks for: \kinds{Paper} by hand, \kinds{CircuitVerse} in the browser, \kinds{VHDL} in a
module of your own.

\textbf{How to check your work.} From Exercise~\ref{c1:ex:2} on, every exercise that asks you to
write a VHDL module ships a self-checking testbench under \file{lectures/L01/exercises} in the
companion repository. Write your module in its \file{exercises/<module>} directory, using the entity
name and the \textbf{port order} the exercise specifies, then run it with GHDL; \secref{c2:app:c}
has the three commands.

No FPGA board is needed for any exercise. Your job is to get the VHDL right, and the testbench is
how you confirm it. Short answers to the parts you work on paper are in
Appendix~\ref{app:answers}; reference solutions to the VHDL are published in the companion
repository, under \file{lectures/L01/solutions}. Try each exercise yourself before looking at
either.
\end{aside}

% ------------------------------------------------------------------------------------------
\exerciseset{Truth Tables and Gates}
\label{c1:set:1}

\exercise{A function that ignores an input}{Paper, CircuitVerse}
\label{c1:ex:1}
A function of two inputs \code{A}, \code{B} and one output \code{X} has the truth table below:

\begin{narrowtable}{@{}ccc@{}}
  \thead{A} & \thead{B} & \thead{X} \\
  \midrule
  0 & 0 & 1 \\
  0 & 1 & 0 \\
  1 & 0 & 1 \\
  1 & 1 & 0 \\
\end{narrowtable}

\begin{parts}
  \item Derive the sum-of-products equation for \code{X} directly from the truth table
    (\secref{c1:a2}). Two rows have \code{X = 1}, so you get two AND-terms.
  \item Simplify it algebraically. Factor out what the two terms share, and apply
    \code{A + A' = 1} and \code{1 · B' = B'} from \secref{c1:a2}. You should end up with a single
    literal.
  \item Your answer to \textbf{b)} says something about \code{A}. What? Which single gate from
    \secref{c1:a1}'s table computes the function, and what happens to the \code{A} input when you
    draw it?
  \item Build that single-gate version in CircuitVerse and confirm it reproduces every row of the
    truth table above. Leave the \code{A} input unconnected, or simply leave it out of the drawing;
    \textbf{b)} is the justification for doing so.
\end{parts}

% ------------------------------------------------------------------------------------------
\exerciseset{Entities and Architectures}
\label{c1:set:2}

\exercise{AND and NAND}{VHDL}
\label{c1:ex:2}
Write two small gates as complete VHDL modules, \code{and_gate} and \code{nand_gate}.

Neither gate is the point, and you should not expect either to be difficult. The point is that
this is where an entity and an architecture stop being something you read in \secref{c1:a4} and
become something you write, and that the two halves are worth writing in that order.

Both entities have the same ports:

\begin{narrowtable}{@{}lll@{}}
  \thead{Port} & \thead{Direction} & \thead{Type} \\
  \midrule
  \code{a} & in & \code{std_logic} \\
  \code{b} & in & \code{std_logic} \\
  \code{x} & out & \code{std_logic} \\
\end{narrowtable}

\begin{parts}
  \item Write the \code{entity} declaration for \code{and_gate}, and nothing else. In VHDL a
    module's interface is a thing you can write down and reason about before any implementation
    exists at all, which is the habit this part exists to build. Identify the required ports, the
    direction of each, and the type of each.
  \item Add an \code{architecture} implementing the AND gate with a single concurrent signal
    assignment.
  \item Now write \code{nand_gate}: the same interface with an inverted output, using VHDL's
    \code{nand} operator directly. Do \textbf{not} write \code{not (a and b)}.
\end{parts}
The two are equivalent as Boolean algebra, and on an FPGA they make no difference at all, because
both end up in the same lookup table. It matters on the kind of hardware Exercise~\ref{c1:ex:6} is
about, where NAND is the gate you physically have and everything else has to be built from it.
Writing what you mean, and letting the tools decide what it costs, is the habit; reaching past an
operator the language already gives you is the thing to avoid.

\modulefigure{and_gate}
\modulefigure{nand_gate}

\selfcheck name your entities \code{and_gate} and \code{nand_gate}, each with ports \code{a},
\code{b} (in) and \code{x} (out), declared in that order; their testbenches are in
\file{exercises/and_gate} and \file{exercises/nand_gate}, and each checks all four rows of the
corresponding truth table.

\exercise{Three-input XOR}{Paper, VHDL, CircuitVerse}
\label{c1:ex:3}
An XOR gate has three inputs, \code{a}, \code{b}, \code{c}, and one output, \code{x}. The output is
\code{1} whenever an odd number of the inputs are \code{1} (the 3-input XOR, or parity, function).

\begin{parts}
  \item Write the truth table for \code{x} over all eight combinations of \code{a}, \code{b},
    \code{c}, and confirm it matches the ``odd number of \code{1}s'' description.
  \item Implement the gate in VHDL, as a module of your own with an entity whose ports are three
    \code{std_logic} inputs and one \code{std_logic} output, and an architecture that drives
    \code{x} with a single concurrent assignment (\secref{c1:a4} shows the entity/architecture
    pattern and notes that VHDL's \code{xor} operator works directly on \code{std_logic};
    \secref{c1:a5} shows a complete module end to end).
  \item Build the same gate in CircuitVerse and confirm it reproduces every row of your truth
    table.
\end{parts}
\note[Tip] \code{xor} chains left-to-right, so the whole architecture body is a single line,
\code{x <= a xor b xor c;}, with no internal \code{signal} or intermediate gates needed.

\modulefigure{xor3}

\selfcheck name your VHDL entity \code{xor3}, with ports \code{a}, \code{b}, \code{c} (in) and
\code{x} (out), declared in that order; its testbench is in \file{exercises/xor3}, and
\secref{c1:a5} lists the three commands that run it.

\exercise{Majority gate}{Paper, VHDL, CircuitVerse}
\label{c1:ex:4}
A \textbf{majority gate} has three inputs, \code{a}, \code{b}, \code{c}, and one output, \code{x}.
The output is \code{1} whenever \emph{at least two} of the three inputs are \code{1}.

This is the circuit behind a triple-redundant vote: three sensors report the same measurement, and
the system acts on whatever answer two of them agree on, so a single failed sensor cannot decide
anything on its own.

\begin{parts}
  \item Write the truth table for \code{x} over all eight combinations of \code{a}, \code{b},
    \code{c}.
  \item Derive the sum-of-products equation for \code{x} directly from the truth table
    (\secref{c1:a2}):
    \begin{itemize}
      \item Four rows have \code{x = 1}, so you get four AND-terms.
      \item Each term includes all three inputs, primed where that input is \code{0} in its row.
    \end{itemize}
  \item Implement the gate in VHDL, as a module of your own:
    \begin{itemize}
      \item An entity with three \code{std_logic} inputs and one \code{std_logic} output.
      \item An architecture driving \code{x} with a single concurrent assignment, translating your
        four AND-terms directly (\secref{c1:a4}).
    \end{itemize}
  \item Build the same network in CircuitVerse and confirm it reproduces every row of your truth
    table.
\end{parts}
\note[Tip] Write the equation out on paper before you open either the editor or CircuitVerse. Four
three-input AND-terms plus one four-input OR gate is a lot of wiring to correct after the fact.

\note[Looking ahead] count the gates your equation needs, then convince yourself by inspection that
\code{x = ab + ac + bc} computes the same function with three two-input AND gates instead of four
three-input ones. Sum-of-products gave you a \emph{correct} network, not the \emph{smallest} one;
finding the smaller one systematically, rather than by staring at it, is what Chapter~\ref{c2:ch}
is for.

\modulefigure{majority3}

\selfcheck name your VHDL entity \code{majority3}, with ports \code{a}, \code{b}, \code{c} (in) and
\code{x} (out), declared in that order; its testbench is in \file{exercises/majority3} and checks
all eight rows. Either equation passes it: the testbench checks the function, not which network
you chose.

\exercise{Half adder}{VHDL, CircuitVerse}
\label{c1:ex:5}
Write a \textbf{half adder}: \code{sum} and \code{carry} from two single bits. You know the circuit;
the reason it is here is that it is the first module in the book with \textbf{two outputs}, and that
is a VHDL question rather than a logic one.

\begin{parts}
  \item Implement it, with two concurrent assignments and no internal \code{signal}.

    Two outputs does not mean two entities, two architectures, or a process. Each output port is
    driven by its own concurrent assignment, and the two are not steps that happen in an order:
    they describe two separate pieces of wire, both live all the time. If you find yourself
    reaching for a process to ``return'' two values, that instinct is the software one, and
    \secref{c1:a4}'s ``a concurrent assignment describes a wire'' is the correction.
  \item Build it in CircuitVerse and confirm both outputs, so you have seen the two-output shape as
    a drawing before you meet it as an entity with several \code{out} ports in
    Chapter~\ref{c2:ch}.
\end{parts}
\note[Looking ahead] this is the cell Chapter~\ref{c6:ch}'s counter is built from. The 4-bit adder
there is four of these chained carry-to-carry, and the counter is that adder with its output fed
back to its input, which is why the wraparound needs no logic of its own.

\modulefigure{half_adder}

\selfcheck name your entity \code{half_adder}, with ports \code{a}, \code{b} (in) and \code{sum},
\code{carry} (out), declared in that order; its testbench is in \file{exercises/half_adder} and
checks \code{sum} and \code{carry} separately, so a design that has one right and the other wrong is
told which is which.

% ------------------------------------------------------------------------------------------
\exerciseset{Universal Gates}
\label{c1:set:3}

\exercise{OR from NAND}{Paper, CircuitVerse, VHDL}
\label{c1:ex:6}
\Secref{c1:a1} states that any Boolean function can be built using only \code{NAND} gates. This
exercise makes you prove it to yourself for the three operators you have been using.

\begin{parts}
  \item Derive the \code{OR} case on paper before you build anything, because De Morgan
    (\secref{c1:laws}) gives it to you directly. Start from \code{(A'B')'} and apply the law; you
    should arrive at \code{A + B}. Now read that expression back as a circuit: what is \code{A'B'}
    with an inversion on its output, and how many \code{NAND} gates is the whole thing?
  \item In CircuitVerse, build each of the following from \code{NAND} gates and nothing else,
    confirming each against its truth table in \secref{c1:a1} as you go:
    \begin{itemize}
      \item \code{NOT a}: one NAND gate is enough, if you feed the same signal to both of its
        inputs.
      \item \code{a AND b}: one NAND gate, followed by the \code{NOT} you just built.
      \item \code{a OR b}: the network you derived in \textbf{a)}. Check your gate count against
        what you predicted.
    \end{itemize}
  \item Write the \code{OR} version in VHDL, as a module of your own:
    \begin{itemize}
      \item An entity with two \code{std_logic} inputs and one \code{std_logic} output.
      \item An architecture driving \code{x} with a single concurrent assignment.
      \item Use only VHDL's \code{nand} operator:
      \begin{itemize}
        \item no \code{or}, no \code{and}, no \code{not} anywhere in the expression.
      \end{itemize}
    \end{itemize}
  \item Explain why this matters in practice. Consider that a chip fabrication process has to
    physically implement every gate type it offers, and that \code{NAND} is the cheapest gate to
    build in CMOS.
\end{parts}
\note[Tip] VHDL's \code{nand} operator does not chain the way \code{and} and \code{or} do:
\code{a nand b nand c} is not legal. Parenthesize every NAND explicitly, which is what you want
here anyway, since each pair of parentheses is one physical gate in your CircuitVerse drawing.

\modulefigure{or_from_nand}

\selfcheck name your VHDL entity \code{or_from_nand}, with ports \code{a}, \code{b} (in) and
\code{x} (out), declared in that order; its testbench is in \file{exercises/or_from_nand} and checks
it against all four rows of the \code{OR} truth table.

% ------------------------------------------------------------------------------------------
\exerciseset{Putting It Together}
\label{c1:set:4}

\exercise{The braking assistant}{VHDL, Paper}
\label{c1:ex:7}
Write the braking assistant from \secref{c1:a6} in VHDL.

\Secref{c1:a6} had you derive its equations, count its gates and sketch its network \emph{before}
reading on. This is where you write it yourself, and where a testbench decides whether your reading
of the requirement matches the one in \secref{c1:a6}'s truth table.

The entity must have:

\begin{booktable}{@{}p{7.5em}p{4.8em}p{5.2em}L@{}}
  \thead{Port} & \thead{Direction} & \thead{Type} & \thead{Description} \\
  \midrule
  \code{driver_brake} & in & \code{std_logic} & The driver is on the brake pedal. \\
  \code{sensor} & in & \code{std_logic} & The proximity sensor sees an obstacle. \\
  \code{radar} & in & \code{std_logic} & The radar sees an obstacle. \\
  \code{error} & in & \code{std_logic} & The assistance system is reporting a fault. \\
  \code{engine_brake} & out & \code{std_logic} & Apply the brakes. \\
\end{booktable}

\begin{parts}
  \item Implement it from \textbf{your own} equations rather than from anyone else's. If your
    equations and the truth table disagree, that disagreement is the most useful thing this
    exercise can give you, and the testbench will tell you which way.

    Use two concurrent assignments and one internal \code{signal}:
    \begin{itemize}
      \item The signal holds the assistance system's own decision, before the driver's pedal is
        considered.
      \item This is the first exercise in the book with an internal signal, and it is the reason the
        keyword exists: a name for an intermediate result, so each line still reads like the clause
        of the requirement it came from.
    \end{itemize}
  \item Now break it deliberately, in the specific way \secref{c1:a6} warns about. Feed
    \code{driver_brake} \emph{through} the fault logic instead of letting it bypass:

\begin{vhdlcode}
engine_brake <= (driver_brake or sensor or radar) and (not error);
\end{vhdlcode}

    That is one gate simpler, and it looks reasonable. Run the testbench against it and answer:
    \begin{itemize}
      \item Which input combination does it fail on first, and what is physically happening in the
        car at that moment?
      \item How many of the sixteen rows does it get wrong, and what do they all have in common?
      \item The failing rows are exactly the ones the safety requirement was written for. What does
        that tell you about the value of a truth table you derived from a requirement, as opposed
        to one you read off a circuit you had already built?
    \end{itemize}
\end{parts}

\modulefigure{adas}

\selfcheck name your entity \code{adas}, with inputs \code{driver_brake}, \code{sensor},
\code{radar}, \code{error} and output \code{engine_brake}, declared in that order; its testbench is
in \file{exercises/adas} and checks all sixteen rows of \secref{c1:a6}'s truth table.
